\section*{Sissejuhatus}\label{sissejuhatus} \addcontentsline{toc}{section}{Sissejuhatus}

Kompilaatorite õpetamisel on oluline näidata, kuidas kõrgtaseme programmeerimiskeeltes kirjutatud lähtekood tõlgitakse madalama taseme käskudeks ning kuidas need käsud lõpuks täidetakse. Hea viis selle protsessi illustreerimiseks on abstraktsete masinate kasutamine. \textsc{CMa} on abstraktne masin, mis modelleerib kompileeritud koodi täitmist pinupõhises keskkonnas. Mudel võimaldab näha, kuidas kompileeritud lähtekoodist saavad konkreetsed käsud, mida virtuaalmasin interpreteerib ning kuidas toimuvad selle tulemusena reaalsed muudatused mälus ja registrites, analoogselt päris riistvara tööga.

Tartu Ülikooli arvutiteaduse instituudis on loodud \textsc{CMa} mudeli implementatsioon Javas, mida kasutatakse ``Automaadid, keeled ja translaatorid'' (edaspidi AKT) kursusel õpetusliku tööriistana. Aine üks eesmärkidest on näidata, kuidas transleerida programmi abstraktne süntaksipuu virtuaalmasinale arusaadavaks käskudeks\footnote{\url{https://ois2.ut.ee/\#/courses/LTAT.03.006/details}}, mille sihtkeele tõlgendamist tööriist toetab. Praegune implementatsioon toetab põhilisi \textsc{CMa} käske, kuid ei võimalda funktsioonikutsete täitmist. Funktsioonikutse\footnote{\url{https://akit.cyber.ee/term/16364-function-call}} (ingl \emph{function call}) on imperatiivsetes programmeerimiskeeltes keskne mehhanism, mille abil antakse juhtimine üle alamprogrammile, edastatakse argumendid ning taastatakse pärast täitmist eelnev täitmiskontekst. Selle toe puudumine piirab simulaatori kasutusvõimalusi nii õppe- kui ka uurimistöös, sest ilma funktsioonikutseteta ei ole võimalik veenvalt modelleerida protseduurilise programmi loomulikku juhtvoogu.

Käesoleva bakalaureusetöö eesmärk oli laiendada \textsc{CMa} mudeli Java-implementatsiooni, lisades funktsioonikutsete toe. Oluline oli tagada, et täitmiskeskkond, sealhulgas pinu\footnote{\url{https://akit.cyber.ee/term/10638-stack-2}} (ingl \emph{stack}), jääb igas olukorras korrektseks, kaasa arvatud pesastatud\footnote{\url{https://akit.cyber.ee/term/5520-nest}} (ingl \emph{nested}) ja rekursiivsete väljakutsete korral. Töö fookus piirdus üksnes imperatiivsete programmeerimiskeelte käsitlemisega. Kõik laiendused ja analüüsid viidi läbi selle kontekstis, lähtudes Wilhelm ja Seidli käsitlusest~\cite{wilhelm_imperative_2010} ning teiste programmeerimisparadigmade käsitlus jäi töö skoobist välja.

Esmalt antakse ülevaade abstraktsete masinate rollist kompilaatorites, kirjeldatakse \textsc{CMa} arhitektuuri ning võrreldakse seda teiste levinud virtuaalmasinatega, nagu \textsc{JVM} ja \textsc{LLVM}. Seejärel kirjeldatakse olemasolevat \textsc{CMa} koodibaasi, analüüsitakse funktsioonikutsete toe puudumist ning sõnastatakse laienduse nõuded. Järgnevalt käsitletakse funktsioonikutsete realiseerimist, keskkonna korrektsuse tagamist ning näidatakse tehnilisi lahendusi. Pärast seda esitatakse testimise tulemused ja võrreldakse lahendust eelneva seisuga. Lõpuks rõhutatakse olulisemaid tulemusi ja tuuakse välja võimalikud edasiarenduse suunad.

Töö tulemusena valmis laiendatud \textsc{CMa} simulaator, mis toetab 13 uut masinakäsku ning on võimeline täitma protseduurilisi programme, sh rekursiivseid väljakutseid. Lahenduse korrektsus kinnitati 46 JUnit testiga ning kokku lisati 1251 rida uut koodi.

Töö sõnastuse ja struktuuri viimistlemisel kasutati abivahendina tehisintellekti. Teksti toimetamiseks, keelelise selguse parandamiseks ning tehniliste mõistete ühtlustamiseks rakendati OpenAI poolt loodud keelemudelit GPT-5.5\footnote{\url{https://developers.openai.com/api/docs/models/gpt-5.5}}. Mudelit kasutati eelkõige soovituste ja alternatiivsete sõnastuste pakkumiseks, kuid lõplikud otsused ja valikud tegi autor ise, tagades töö sisulise korrektsuse ja vastavuse akadeemilistele nõuetele.
